\documentclass[11pt,preprint]{elsarticle}

\usepackage{lmodern}
%%%% My spacing
\usepackage{setspace}
\setstretch{1.2}
\DeclareMathSizes{12}{14}{10}{10}

% Wrap around which gives all figures included the [H] command, or places it "here". This can be tedious to code in Rmarkdown.
\usepackage{float}
\let\origfigure\figure
\let\endorigfigure\endfigure
\renewenvironment{figure}[1][2] {
    \expandafter\origfigure\expandafter[H]
} {
    \endorigfigure
}

\let\origtable\table
\let\endorigtable\endtable
\renewenvironment{table}[1][2] {
    \expandafter\origtable\expandafter[H]
} {
    \endorigtable
}


\usepackage{ifxetex,ifluatex}
\usepackage{fixltx2e} % provides \textsubscript
\ifnum 0\ifxetex 1\fi\ifluatex 1\fi=0 % if pdftex
  \usepackage[T1]{fontenc}
  \usepackage[utf8]{inputenc}
\else % if luatex or xelatex
  \ifxetex
    \usepackage{mathspec}
    \usepackage{xltxtra,xunicode}
  \else
    \usepackage{fontspec}
  \fi
  \defaultfontfeatures{Mapping=tex-text,Scale=MatchLowercase}
  \newcommand{\euro}{€}
\fi

\usepackage{amssymb, amsmath, amsthm, amsfonts}

\def\bibsection{\section*{References}} %%% Make "References" appear before bibliography


\usepackage[numbers]{natbib}

\usepackage{longtable}
\usepackage[margin=2.3cm,bottom=2cm,top=2.5cm, includefoot]{geometry}
\usepackage{fancyhdr}
\usepackage[bottom, hang, flushmargin]{footmisc}
\usepackage{graphicx}
\numberwithin{equation}{section}
\numberwithin{figure}{section}
\numberwithin{table}{section}
\setlength{\parindent}{0cm}
\setlength{\parskip}{1.3ex plus 0.5ex minus 0.3ex}
\usepackage{textcomp}
\renewcommand{\headrulewidth}{0.2pt}
\renewcommand{\footrulewidth}{0.3pt}

\usepackage{array}
\newcolumntype{x}[1]{>{\centering\arraybackslash\hspace{0pt}}p{#1}}

%%%%  Remove the "preprint submitted to" part. Don't worry about this either, it just looks better without it:
\makeatletter
\def\ps@pprintTitle{%
  \let\@oddhead\@empty
  \let\@evenhead\@empty
  \let\@oddfoot\@empty
  \let\@evenfoot\@oddfoot
}
\makeatother

 \def\tightlist{} % This allows for subbullets!

\usepackage{hyperref}
\hypersetup{breaklinks=true,
            bookmarks=true,
            colorlinks=true,
            citecolor=blue,
            urlcolor=blue,
            linkcolor=blue,
            pdfborder={0 0 0}}


% The following packages allow huxtable to work:
\usepackage{siunitx}
\usepackage{multirow}
\usepackage{hhline}
\usepackage{calc}
\usepackage{tabularx}
\usepackage{booktabs}
\usepackage{caption}


\newenvironment{columns}[1][]{}{}

\newenvironment{column}[1]{\begin{minipage}{#1}\ignorespaces}{%
\end{minipage}
\ifhmode\unskip\fi
\aftergroup\useignorespacesandallpars}

\def\useignorespacesandallpars#1\ignorespaces\fi{%
#1\fi\ignorespacesandallpars}

\makeatletter
\def\ignorespacesandallpars{%
  \@ifnextchar\par
    {\expandafter\ignorespacesandallpars\@gobble}%
    {}%
}
\makeatother


% definitions for citeproc citations
\NewDocumentCommand\citeproctext{}{}
\NewDocumentCommand\citeproc{mm}{%
\href{\#cite.\detokenize{#1}}{#2}\nocite{#1}}

\makeatletter
% allow citations to break across lines
\let\@cite@ofmt\@firstofone
% avoid brackets around text for \cite:
\def\@biblabel#1{}
\def\@cite#1#2{{#1\if@tempswa , #2\fi}}
\makeatother
\newlength{\cslhangindent}
\setlength{\cslhangindent}{1.5em}
\newlength{\csllabelwidth}
\setlength{\csllabelwidth}{3em}
\newenvironment{CSLReferences}[2] % #1 hanging-indent, #2 entry-spacing
{\begin{list}{}{%
	\setlength{\itemindent}{0pt}
	\setlength{\leftmargin}{0pt}
	\setlength{\parsep}{0pt}
	% turn on hanging indent if param 1 is 1
	\ifodd #1
	\setlength{\leftmargin}{\cslhangindent}
	\setlength{\itemindent}{-1\cslhangindent}
	\fi
	% set entry spacing
	\setlength{\itemsep}{#2\baselineskip}}}
{\end{list}}

\usepackage{calc}
\newcommand{\CSLBlock}[1]{\hfill\break\parbox[t]{\linewidth}{\strut\ignorespaces#1\strut}}
\newcommand{\CSLLeftMargin}[1]{\parbox[t]{\csllabelwidth}{\strut#1\strut}}
\newcommand{\CSLRightInline}[1]{\parbox[t]{\linewidth - \csllabelwidth}{\strut#1\strut}}
\newcommand{\CSLIndent}[1]{\hspace{\cslhangindent}#1}


\urlstyle{same}  % don't use monospace font for urls
\setlength{\parindent}{0pt}
\setlength{\parskip}{6pt plus 2pt minus 1pt}
\setlength{\emergencystretch}{3em}  % prevent overfull lines
\setcounter{secnumdepth}{5}

%%% Use protect on footnotes to avoid problems with footnotes in titles
\let\rmarkdownfootnote\footnote%
\def\footnote{\protect\rmarkdownfootnote}
\IfFileExists{upquote.sty}{\usepackage{upquote}}{}

%%% Include extra packages specified by user
\providecommand{\pandocbounded}[1]{#1}

%%% Hard setting column skips for reports - this ensures greater consistency and control over the length settings in the document.
%% page layout
%% paragraphs
\setlength{\baselineskip}{12pt plus 0pt minus 0pt}
\setlength{\parskip}{12pt plus 0pt minus 0pt}
\setlength{\parindent}{0pt plus 0pt minus 0pt}
%% floats
\setlength{\floatsep}{12pt plus 0 pt minus 0pt}
\setlength{\textfloatsep}{20pt plus 0pt minus 0pt}
\setlength{\intextsep}{14pt plus 0pt minus 0pt}
\setlength{\dbltextfloatsep}{20pt plus 0pt minus 0pt}
\setlength{\dblfloatsep}{14pt plus 0pt minus 0pt}
%% maths
\setlength{\abovedisplayskip}{12pt plus 0pt minus 0pt}
\setlength{\belowdisplayskip}{12pt plus 0pt minus 0pt}
%% lists
\setlength{\topsep}{10pt plus 0pt minus 0pt}
\setlength{\partopsep}{3pt plus 0pt minus 0pt}
\setlength{\itemsep}{5pt plus 0pt minus 0pt}
\setlength{\labelsep}{8mm plus 0mm minus 0mm}
\setlength{\parsep}{\the\parskip}
\setlength{\listparindent}{\the\parindent}
%% verbatim
\setlength{\fboxsep}{5pt plus 0pt minus 0pt}



\begin{document}



\begin{frontmatter}  %

\title{Comparing the JSE Large, Mid and Small Cap indices}

% Set to FALSE if wanting to remove title (for submission)




\author[Add1]{24693197}
\ead{24693197@sun.ac.za}





\address[Add1]{Stellenbosch University, Stellenbosch, South Africa}



\vspace{1cm}





\vspace{0.5cm}

\end{frontmatter}

\setcounter{footnote}{0}



%________________________
% Header and Footers
%%%%%%%%%%%%%%%%%%%%%%%%%%%%%%%%%
\pagestyle{fancy}
\chead{}
\rhead{}
\lfoot{}
\rfoot{\footnotesize Page \thepage}
\lhead{}
%\rfoot{\footnotesize Page \thepage } % "e.g. Page 2"
\cfoot{}

%\setlength\headheight{30pt}
%%%%%%%%%%%%%%%%%%%%%%%%%%%%%%%%%
%________________________

\headsep 35pt % So that header does not go over title




\section{\texorpdfstring{Introduction
\label{Introduction}}{Introduction }}\label{introduction}

This report compares the traits of the three local size indices, namely
the large cap, mid cap and small cap indices. The Capped SWIX is used as
the benchmark because it represents the broader local equity market. The
report first looks at what each index holds (its sector exposure and
stock concentration). It then compares the risk of each index, using the
realised volatility of its constituents and the behaviour of each index
in high and low volatility periods. Lastly, it compares the returns of
the indices and how much each differs from the benchmark, using active
share. Small cap stocks are usually expected to be more volatile than
large cap stocks. This report will therefore also look at whether this
holds for the local size indices.

\section{Data and appraoch}\label{data-and-appraoch}

The report uses the monthly constituents of each index from 2007, along
with daily stock and index returns from 2015. Therefore, the volatility
and return results cover a shorter period than the sector and
concentration results.

The following techniques are used:

\begin{itemize}
\tightlist
\item
  Concentration refers to how much of an index is held in its largest
  stocks. A market capitalisation weighted index can become too
  concentrated in a few large stocks. It is measured as the combined
  weight of the ten largest stocks and ten was chosen as the large cap
  index holds only around forty stocks.
\item
  Realised volatility refers to the annualised standard deviation of
  daily returns. Standard deviation is mostly used as a means of
  estimating risk. For each stock, it is calculated over a rolling one
  year window. The volatility of an index's constituents is then the
  weighted average of these values, using the index weights at each
  month end.
\item
  A drawdown refers to the fall from the previous peak. For each stock,
  the largest drawdown within each calendar year is calculated. The
  median across the constituents of each index is then compared.
\item
  Stratification refers to studying the statistical properties of a
  series during different periods. The periods are stratified according
  to the monthly volatility of the Capped SWIX. Months are split into
  high volatility (above the 85th percentile), low volatility (below the
  25th percentile) and other months. The performance of each index is
  compared across these periods.
\item
  Principal component analysis (PCA) aims to reduce the dimensionality
  of the sources of variation in a set of returns. It reorganises the
  returns into independent components that explain (in decreasing order)
  the variation in the data. For each index and calendar year, PCA is
  applied to the correlation matrix of the constituents' daily log
  returns. The correlation matrix was chosen as it lets all stocks enter
  on an equal volatility basis. The proportion of variation explained by
  the first component is used as a measure of how much the stocks in an
  index move together.
\item
  Rolling returns refer to the annualised return over a moving three
  year window. Cumulative returns can be misleading because early
  outperformance can greatly skew later performance. Rolling returns
  avoid this and show the experience of an investor who held the index
  for three years.
\item
  Active share refers to the share of an index's holdings that differ
  from the benchmark. It ranges from zero (the same stocks and weights
  as the Capped SWIX) to one (no stocks in common). It is calculated as
\end{itemize}

\begin{align}
AS_t = \frac{1}{2} \sum_{i=1}^{N} \left| w_{i,t}^{Index} - w_{i,t}^{SWIX} \right| \label{eq1}
\end{align}

where a stock not held by one of the two indices has a weight of zero.

\section{\texorpdfstring{Sector Exposure and Concentration
\label{Comp}}{Sector Exposure and Concentration }}\label{sector-exposure-and-concentration}

This section will look at what each index holds (sector exposure) and
each index's concentration.

\begin{figure}
\centering
\pandocbounded{\includegraphics[keepaspectratio,alt={Sector exposure by index }]{Question2_files/Question2_files/figure-latex/sector-1.png}}
\caption{Sector exposure by index \label{fig:sector}}
\end{figure}

Figure \ref{fig:sector} shows the weight of each sector in each index
over time. Industrials is the largest sector in the Capped SWIX over the
whole sample and the second largest is resources. The large cap index
started by holding close to 50\% of its weight in resources. This share
fell to its lowest point around 2016 and has increased again since. The
large cap index also had an increased share of financials held from
around 2015. The mid cap index was mainly made up of industrials at the
start of the sample, but this share has declined as financials and
resources have grown. The small cap index has the largest exposure to
property of all the indices. A large part of the small cap index has no
sector label early in the sample, but this share has decreased since
2018. This figure shows that that each index carries a different sector
risk.

\begin{figure}
\centering
\pandocbounded{\includegraphics[keepaspectratio,alt={Weight of the ten largest stocks in each index }]{Question2_files/Question2_files/figure-latex/conc-1.png}}
\caption{Weight of the ten largest stocks in each index
\label{fig:conc}}
\end{figure}

Figure \ref{fig:conc} shows the combined weight of the ten largest
stocks in each index. A market capitalisation weighted index can become
concentrated in a few large stocks. The large cap index is the most
concentrated throughout the sample, although its top ten weight fell
after 2023 and in recent years it is only slightly higher than the
others. The mid cap index has become more concentrated since 2020 and is
close to the large cap index by the end of the sample. The small cap
index is the least concentrated for most of the sample. Until around
2020, the Capped SWIX lies between the large cap index and the smaller
indexes. Therefore, an investor in the large cap index is more exposed
to the performance of a few stocks than an investor in the small cap
index.

\section{\texorpdfstring{Volatility
\label{Vol}}{Volatility }}\label{volatility}

This section will look at the risk of each index, using a few different
techniques to evaluate and discuss it.

\begin{figure}
\centering
\pandocbounded{\includegraphics[keepaspectratio,alt={Rolling one year realised volatility of index constituents }]{Question2_files/Question2_files/figure-latex/cons_vol-1.png}}
\caption{Rolling one year realised volatility of index constituents
\label{fig:consvol}}
\end{figure}

Standard deviation is mostly used as a means of estimating risk. Figure
\ref{fig:consvol} shows the weighted average realised volatility of the
constituents of each index. The volatility of all four indices increased
in 2020 and stayed high for a year (as the 2020 returns remain in the
one year window)). The small cap constituents are the most volatile for
most of 2016 to 2022. The large cap index and the Capped SWIX follow
almost the same path. From 2022 to 2024 the mid cap constituents are the
least volatile. From 2025 the small cap constituents are the least
volatile.

\begin{figure}
\centering
\pandocbounded{\includegraphics[keepaspectratio,alt={Median largest peak-to-trough fall of index constituents within each year }]{Question2_files/Question2_files/figure-latex/year_dd-1.png}}
\caption{Median largest peak-to-trough fall of index constituents within
each year \label{fig:yeardd}}
\end{figure}

Figure \ref{fig:yeardd} shows the median of each stock's largest fall
from its peak within each year. In most years, the typical stock in each
index fell by between a fifth and a third from its peak. These declines
are very similar across the size indices. The typical large cap stock
falls almost as much as the typical small cap stock in most years. The
largest falls for all indices occurred in 2020, which is to be expected.
This suggests that, the large cap index is about as risky as the small
cap index.

\begin{table}[H]
\centering
\begin{tabular}{llrr}
  \hline
Enviro & Index & SD & Ret \\ 
  \hline
High volatility & CappedSWIX & 0.309 & -0.067 \\ 
  High volatility & LargeCaps & 0.307 & -0.030 \\ 
  High volatility & MidCaps & 0.293 & -0.081 \\ 
  High volatility & SmallCaps & 0.254 & -0.092 \\ 
  Low volatility & CappedSWIX & 0.101 & 0.157 \\ 
  Low volatility & LargeCaps & 0.114 & 0.211 \\ 
  Low volatility & MidCaps & 0.114 & 0.095 \\ 
  Low volatility & SmallCaps & 0.082 & 0.107 \\ 
  Other & CappedSWIX & 0.159 & 0.116 \\ 
  Other & LargeCaps & 0.173 & 0.123 \\ 
  Other & MidCaps & 0.154 & 0.123 \\ 
  Other & SmallCaps & 0.118 & 0.158 \\ 
   \hline
\end{tabular}
\caption{Annualised return and SD by Capped SWIX volatility environment \label{tab:strat}} 
\end{table}

Table \ref{tab:strat} shows the annualised return and standard deviation
of each index in high volatility, low volatility and other months. All
four indices have a negative average return in high volatility months.
In these months the large cap index falls the least and the small cap
index falls the most. In low volatility months the large cap index has
the highest return. In other months the small cap index has the highest
return. The small cap index has the lowest standard deviation in every
group. This contrasts with Figure \ref{fig:consvol}, where its
constituents are the most volatile for much of the sample.

\begin{figure}
\centering
\pandocbounded{\includegraphics[keepaspectratio,alt={Proportion of constituent return variation explained by the first principal component }]{Question2_files/Question2_files/figure-latex/pca-1.png}}
\caption{Proportion of constituent return variation explained by the
first principal component \label{fig:pca}}
\end{figure}

Figure \ref{fig:pca} shows the proportion of the variation in each
index's constituent returns that is explained by the first principal
component. The first principal component can be interpreted as a common
market effect. Threfore, a higher proportion implies a higher level of
commonality in returns between the stocks in an index. This leaves less
room for diversification within the index. The large cap index has the
highest proportion in every year and the small cap index has the lowest.
The proportion increased for all indices in 2020 because the comovement
between stocks increased during this period. This helps explain the
earlier results. Although the small cap constituents are volatile, they
show less comovement. Therefore, their movements partly offset each
other and the small cap index has the lowest volatility.

\section{\texorpdfstring{Returns \label{Ret}}{Returns }}\label{returns}

This section will look at the returns of each index.

\begin{figure}
\centering
\pandocbounded{\includegraphics[keepaspectratio,alt={Rolling three year annualised returns }]{Question2_files/Question2_files/figure-latex/roll_rets-1.png}}
\caption{Rolling three year annualised returns \label{fig:rollrets}}
\end{figure}

Figure \ref{fig:rollrets} shows the rolling three year annualised return
of each index. The small cap index has the widest range of rolling
returns. It had the lowest rolling returns around 2020 and the highest
from 2022 to 2024. The large cap index had the highest rolling returns
for most of the period before 2022, where it surpassed by the small cap
index. The mid cap index is at or below the Capped SWIX for most of the
period. By the end of the sample, all four indices have similar rolling
returns. Therefore, the index that performed best depends on the period
in which an investor held it.

\begin{table}[H]
\centering
\begingroup\footnotesize
\begin{tabular}{llrrrr}
  \hline
Period & Info & CappedSWIX & LargeCaps & MidCaps & SmallCaps \\ 
  \hline
Last 3 Years & Cum Returns & 0.696 & 0.678 & 0.554 & 0.840 \\ 
  Last 3 Years & Returns (Ann.) & 0.193 & 0.188 & 0.158 & 0.225 \\ 
  Last 3 Years & Returns Excess (Ann.) & 0.000 & -0.004 & -0.029 & 0.027 \\ 
  Last 3 Years & SD (Ann.) & 0.128 & 0.139 & 0.136 & 0.110 \\ 
  Last 3 Years & Adj. Sharpe Ratio & 0.853 & 0.946 & 0.855 & 1.023 \\ 
  Last 3 Years & Pain Index & 0.026 & 0.028 & 0.036 & 0.017 \\ 
  Last 3 Years & Avg DD & 0.047 & 0.067 & 0.063 & 0.051 \\ 
  Last 3 Years & Tracking Error & 0.000 & 0.022 & 0.076 & 0.118 \\ 
  Last 3 Years & Information Ratio &  & -0.200 & -0.452 & 0.277 \\ 
  Last 3 Years & Beta & 1.000 & 1.074 & 0.887 & 0.448 \\ 
  Last 3 Years & Beta Bull & 1.000 & 1.126 & 0.561 & 0.194 \\ 
  Last 3 Years & Beta Bear & 1.000 & 1.060 & 0.773 & 0.685 \\ 
  Last 3 Years & Up-Down Ratio & 0.000 & 0.600 & 0.440 & 0.400 \\ 
  Last 3 Years & Modified CVaR & -0.083 & -0.086 & -0.089 & -0.061 \\ 
  Last 5 Years & Cum Returns & 0.919 & 1.079 & 0.635 & 1.204 \\ 
  Last 5 Years & Returns (Ann.) & 0.139 & 0.158 & 0.103 & 0.171 \\ 
  Last 5 Years & Returns Excess (Ann.) & 0.000 & 0.016 & -0.032 & 0.028 \\ 
  Last 5 Years & SD (Ann.) & 0.137 & 0.154 & 0.142 & 0.113 \\ 
  Last 5 Years & Adj. Sharpe Ratio & 0.918 & 0.986 & 0.694 & 1.151 \\ 
  Last 5 Years & Pain Index & 0.032 & 0.037 & 0.040 & 0.016 \\ 
  Last 5 Years & Avg DD & 0.069 & 0.075 & 0.061 & 0.044 \\ 
  Last 5 Years & Tracking Error & 0.000 & 0.044 & 0.079 & 0.115 \\ 
  Last 5 Years & Information Ratio &  & 0.420 & -0.455 & 0.279 \\ 
  Last 5 Years & Beta & 1.000 & 1.085 & 0.872 & 0.490 \\ 
  Last 5 Years & Beta Bull & 1.000 & 1.244 & 0.527 & 0.253 \\ 
  Last 5 Years & Beta Bear & 1.000 & 1.019 & 0.900 & 0.789 \\ 
  Last 5 Years & Up-Down Ratio & 0.000 & 0.590 & 0.436 & 0.359 \\ 
  Last 5 Years & Modified CVaR & -0.079 & -0.084 & -0.088 & -0.063 \\ 
   \hline
\end{tabular}
\endgroup
\caption{Risk and return statistics of the size indices relative to the Capped SWIX \label{tab:stats}} 
\end{table}

Table \ref{tab:stats} reports a table of statistics for financial ratios
over the last three and five years, relative to the Capped SWIX. As the
Capped SWIX is the benchmark, its relative measures (such as tracking
error and beta) are not meaningful. The small cap index has the highest
annualised return and the lowest standard deviation over both periods.
Therefore, it also has the highest adjusted Sharpe ratio and the
smallest modified CVaR. Its beta is the lowest and its tracking error
the highest. This suggests that it moves the least with the benchmark.
Its beta in bull-market periods is also much lower than its beta in
bear-market periods. Therefore, the small cap index follows the
benchmark less when the market rises than when it falls. The large cap
index has a beta above one and the lowest tracking error of the three
size indexes. The mid cap index has the lowest return over both periods
and a negative information ratio. This suggests that the small cap index
offered the best return for its level of risk over both periods, and
that the mid cap index offered the lowest return. Although it displayed
signs of slightly higher market risk, the large cap index behaved most
like the benchmark.

\section{\texorpdfstring{Active Share
\label{AS}}{Active Share }}\label{active-share}

This section will look at how different each size index is from the
benchmark.

\begin{figure}
\centering
\pandocbounded{\includegraphics[keepaspectratio,alt={Active share of each size index relative to the Capped SWIX }]{Question2_files/Question2_files/figure-latex/active_share-1.png}}
\caption{Active share of each size index relative to the Capped SWIX
\label{fig:as}}
\end{figure}

Figure \ref{fig:as} shows the active share of each size index relative
to the Capped SWIX. The large cap index has the lowest active share.
This suggests that its stocks make up a large part of the Capped SWIX.
The small cap index has an active share close to one through the whole
sample period. The mid cap index lies between the two, but generally
closer to the small cap index. The active share of the large cap index
decreased in 2023 and has remained at this lower level since. The active
share of the mid cap index slightly declined after 2020 and has
increased again since 2025. Therefore, the small cap index gives an
investor the most different set of stocks from the benchmark, while the
large cap index gives the least.

\section{\texorpdfstring{Conclusion
\label{Conc}}{Conclusion }}\label{conclusion}

This report compared the large, mid and small cap indices with the
Capped SWIX. Each index carries a different sector exposure. The large
cap index is the most concentrated, its stocks show the most comovement
and it is the most similar to the benchmark. The mid cap index has had
the lowest returns in recent years. The small cap constituents are the
most volatile for much of the sample. However, they show the least signs
of comovement. Therefore, the small cap index has the lowest volatility
and the highest risk-adjusted return over the last three and five years.
It is also the most different from the benchmark,but it had the largest
average losses in high volatility months. Contrary to what is widely
expected, the results do not support the view that small caps are the
most volatile part of the local market. The large cap index carries
similar or higher risk than the small cap index over large parts of the
sample period. It is evident that the size indices offer investors
different risk and return profiles. In particular, the small cap index
offers diversification away from the benchmark.

\bibliography{Tex/ref}





\end{document}
